% =============================================================================
%  experience.tex  —  Typesetting Tips & Best Practices
%
%  This section documents real-world lessons learned from using this template
%  across dozens of XMUM assignments over the course of a year.
% =============================================================================

\section{Best Practices}

The tips in this section come from a year of using this template across
XMUM assignments of all sizes — from 10-page weekly reports to a 500-page
community service final report.
Following these practices will save you time and help you avoid common pitfalls.

% -----------------------------------------------------------------------------
\subsection{Modular File Structure (Strongly Recommended)}
% -----------------------------------------------------------------------------

Rather than writing everything inside a single \texttt{main.tex}, split your
content into one \texttt{.tex} file per section and use \verb|\input{}| to
include them:

\begin{verbatim}
main.tex
sections/
    01_introduction.tex
    02_methodology.tex
    03_results.tex
    04_discussion.tex
    05_conclusion.tex
\end{verbatim}

This approach has several advantages:

\begin{itemize}
    \item \textbf{AI-assisted writing is more effective.}
          AI tools (ChatGPT, Claude, Gemini, etc.) have limited context
          windows. Feeding one section at a time means the AI can read your
          entire section without truncation, resulting in higher-quality
          suggestions and edits.
    \item \textbf{Easier collaboration.} Teammates can work on different
          section files simultaneously without merge conflicts.
    \item \textbf{Faster debugging.} When a compilation error occurs, the
          error message will point to the specific section file rather than
          a line number in a 1000-line \texttt{main.tex}.
\end{itemize}

% -----------------------------------------------------------------------------
\subsection{Fast Compilation for Image-Heavy Documents}
% -----------------------------------------------------------------------------

If your document contains many figures, compilation can become slow because
XeLaTeX re-processes every image on every run.

\textbf{Solution:} Use the \texttt{draft} class option during writing.
In draft mode, figures are replaced by placeholder boxes, reducing compile
time from minutes to seconds:

\begin{verbatim}
\documentclass[draft]{xum_review}   % Fast: figures shown as boxes
\documentclass{xum_review}          % Normal: figures rendered fully
\end{verbatim}

Switch back to the normal option only for your final compilation.

% -----------------------------------------------------------------------------
\subsection{Compiling for Turnitin Submission}
% -----------------------------------------------------------------------------

\textbf{Problem:} Turnitin scans every page of your PDF, including headers
and footers. Because this template places \texttt{XIAMEN UNIVERSITY MALAYSIA}
in the header and \texttt{Page X of Y} in the footer of \emph{every} page,
Turnitin may flag these repeated strings as similarity matches, artificially
inflating your similarity score.

\textbf{Solution:} This template provides a dedicated \texttt{[turnitin]}
class option that disables all headers and footers in one step.
When you are ready to generate the PDF for Turnitin submission, change the
very first line of \texttt{main.tex}:

\begin{verbatim}
% Normal version (for your own records and lecturer submission):
\documentclass{xum_review}

% Turnitin version (for upload to the originality checker):
\documentclass[turnitin]{xum_review}
\end{verbatim}

Compile the Turnitin version, upload the resulting PDF, then switch back
to the normal version for your final lecturer submission.

\begin{itemize}
    \item This is a \textbf{one-word change} — no need to comment out
          multiple lines of header/footer configuration.
    \item The body content, margins, and all formatting remain identical
          between the two versions.
    \item Your lecturer's copy (with the XMUM header) and your Turnitin
          copy are compiled from the same source — no risk of accidentally
          submitting different versions.
\end{itemize}
